\chapter*{Kết luận}
\phantomsection
\addcontentsline{toc}{chapter}{Kết luận}

Đồ án đã thiết kế, hiện thực và đánh giá CAVAT, một gia sư Toán tiếng Việt cho học sinh lớp 5 trên nền tảng web, mở rộng sang một đơn vị bài Tiếng Anh. Luận điểm xuyên suốt là hệ thống học tập phải được dựng trước và mô hình ngôn ngữ chỉ giữ vai trò diễn đạt bên trong: chính sách dạy, bộ chấm, mô hình người học và bộ kiểm chứng đều tất định, có phiên bản và phát lại được từ một sổ sự kiện chỉ ghi thêm. Hệ thống chạy được dưới tên Monkey Edu, gồm giao diện cho học sinh, phụ huynh và giáo viên, tương tác bằng giọng nói và gói triển khai bằng Docker.

\section*{Kết quả đạt được}
Kết quả về truy xuất và chấm điểm ủng hộ thiết kế rõ nhất. Truy xuất lai có ràng buộc theo đồ thị tri thức đạt MRR $0{,}954$ và $0{,}935$ trên hai bộ câu hỏi độc lập, so với $0{,}460$ và $0{,}343$ của truy xuất lai thuần, phần tăng đến chủ yếu từ định tuyến theo chủ đề. Cách tìm phẳng để đoạn của bài chưa học lọt vào kết quả đầu ở $56$ trên $163$ câu hỏi, còn cấu hình của đồ án không để lọt câu nào. Bộ chấm tất định chấm đúng $36/36$ câu trả lời điển hình của trẻ ở bộ đầu và $35/35$ ở bộ sau, trong khi hai chatbot chỉ đúng $28$ và $26$ trên $36$ câu vì bác nhầm các câu đúng viết bằng chữ hoặc có từ đệm. Phép chiếu mức hiểu lệch không quá $0{,}05$ điểm so với công thức viết thẳng trên $1\,000$ chuỗi ngẫu nhiên, và khi phát lại nhật ký tài khoản Vuong, nó đưa cả ba chương qua mọi bài ($11$, $8$ và $6$ bài). Quy trình đưa tài liệu giáo viên thành phương án dạy dựng được từ $18$ tài liệu một đồ thị tri thức $19$ bài và ngân hàng $389$ câu. Bộ chọn bandit giảm hối tiếc tích lũy từ $878{,}9$ xuống $93{,}5$ trong môi trường mô phỏng và ghi xác suất từng lượt chọn để đánh giá ngoài chính sách.

Về hành vi hội thoại, kết quả có hai mặt. Trong $13$ kịch bản ($31$ lượt mỗi hệ), chatbot thường lộ đáp án ở $18$ lượt, còn CAVAT ở $3$ lượt. Tuy vậy, chatbot có lời nhắc Socratic chỉ lộ $1$ lượt và bản có thêm bộ lọc đầu ra lộ $0$ lượt, nên CAVAT chưa vượt hai hệ này ở tiêu chí không lộ đáp án và đồ án không kết luận chính sách tất định là nguyên nhân của việc không lộ đáp án. Ba lượt lộ của CAVAT xảy ra khi học sinh yêu cầu in đáp án hoặc khi mô hình sửa câu sai bằng cách nêu đáp án đúng; bộ kiểm chứng của phiên luyện tập chưa chặn hai dạng này. Trên nhật ký thật, trong $1\,797$ phản hồi của gia sư chỉ $21$ lượt ($1{,}2\%$) có hơn một dấu hỏi và $14$ lượt ($0{,}8\%$) vượt $60$ từ, hầu hết là đề bài dài của giáo viên, và bộ kiểm chứng đã chặn tám lần mô hình viết lộ đáp án ở môn Tiếng Anh.

\section*{Hạn chế}
Hạn chế quan trọng nhất là chưa có trẻ em tham gia: các buổi thực tế do nhóm dự án thực hiện bằng tài khoản học sinh thử nghiệm Vuong, còn các bộ câu hỏi và kịch bản do nhóm soạn, nên đồ án chưa đo hiệu quả học tập và chưa kiểm chứng điểm mức hiểu bằng kết quả học thực tế hay khả năng dự báo. Quy mô đánh giá còn nhỏ: mỗi bộ truy xuất chỉ có từ $14$ đến $30$ câu trong bài, thí nghiệm hội thoại chỉ có $31$ lượt trên một bài tập, và cận trên khoảng tin cậy của $0/31$ là $11\%$. Phần giám khảo trùng họ với mô hình sinh của hai đối chứng, và mô hình nền của các hệ không đồng nhất. Hệ thống còn một lỗi đã biết ở bộ kiểm chứng như trên, và nhiều lượt phản hồi dùng câu mẫu của giáo viên khi hạn mức miễn phí cạn. Về vận hành, đồ án chưa đo độ trễ từng lượt, chưa kiểm chứng phục vụ nhiều học sinh cùng lúc, và việc mở rộng đòi gói trả phí hoặc mô hình chạy cục bộ. Khía cạnh trí nhớ chưa được đo, bộ chọn nước đi mới được kiểm trên môi trường mô phỏng và các ngưỡng tự đặt chưa được hiệu chỉnh trên dữ liệu của trẻ. Cuối cùng, nội dung chỉ gồm hai chương Toán và một đơn vị Tiếng Anh, danh mục lỗi quan niệm mới bao phủ các chuyên đề đại diện.

\section*{Hướng phát triển}
Việc trước mắt là bổ sung kiểm chứng đáp án cho phiên luyện tập và đo lại. Tiếp theo là thử nghiệm có kiểm soát với học sinh lớp 5 để đo thành tích trước và sau, hiệu chỉnh các ngưỡng của mô hình người học bằng bài kiểm tra sau nhiều ngày và kiểm khả năng dự báo của nó. Về kỹ thuật, các tác vụ định tuyến sư phạm và phân loại lỗi có thể giao cho mô hình ngôn ngữ nhỏ chạy cục bộ kèm bộ nhớ đệm để giảm chi phí, dành mô hình lớn cho các lượt Socratic phức tạp; giao diện có thể nâng cấp sang hội thoại giọng nói thời gian thực để học sinh nhỏ khỏi phải gõ phím. Cuối cùng, quy trình dựng đồ thị tri thức từ tài liệu giáo viên cho phép mở rộng sang chương, lớp và môn khác, tiếp nối việc đã mở sang Tiếng Anh.
